% ============================================================
% Cheat sheet, argomento 4: dal codice al diagramma, riferimenti fra classi
% ============================================================
\argomento{Dal codice al diagramma: riferimenti fra classi}{Dispensa: capitolo 7. Un campo il cui tipo è una
classe del diagramma diventa un'associazione}[Dal codice: riferimenti fra classi]

\begin{guida}{Come si legge un campo che punta a un'altra classe}
\begin{enumerate}[nosep, leftmargin=*]
  \item Il campo \code{T nome} nella classe \code{A} diventa una linea da \code{A} a \code{T}, con la \textbf{freccia}
        verso \code{T} (da \code{A} si arriva a \code{T}). Se anche \code{T} ha un campo verso \code{A}, la linea è una
        sola e \textbf{senza frecce}.
  \item Il \textbf{nome del campo} è il \textbf{ruolo}: si scrive vicino a \code{T}.
  \item La \textbf{molteplicità vicino a \code{T}} viene dal campo: \code{0..1} se può essere \code{null},
        \code{1} se è sempre impostato (costruttore, \code{final}), \code{*}, \code{1..*}, \code{0..4}\ldots\ se è una
        collezione (argomento 5).
  \item La \textbf{molteplicità vicino a \code{A}} (quanti \code{A} per uno stesso \code{T}) nel codice di \code{A} non
        c'è: viene dal campo inverso, se esiste, o dal testo dell'esercizio. Se nessuno la dice, si lascia vuota.
\end{enumerate}
\end{guida}

\begin{caso}{rif01}{Un riferimento che può mancare}
\begin{lstlisting}[style=javacs]
class Persona {
    private Automobile auto;  // null se non compra
    void compra(Automobile a) { auto = a; }
}
class Automobile {
    private String targa;
}
\end{lstlisting}
\tcblower
\begin{tikzpicture}
  \node[classe1=1.6cm] (p) at (0,0) {Persona};
  \umlclass[2.4cm]{a}{(4.6,-0.12)}{Automobile}{- targa: String}{}
  \draw[navig] (p.east) -- (p.east -| a.west);
  \estO{(p.east -| a.west)}{0..1}{auto}
\end{tikzpicture}
\regola{Freccia da chi ha il campo verso il tipo del campo; ruolo \code{auto}; \code{0..1} perché \code{auto} può
valere \code{null}. \code{Automobile} non ha campi verso \code{Persona}: non è navigabile all'indietro (si può
aggiungere la croce, \vv{navigazione}).}
\end{caso}

\begin{caso}{rif1}{Un riferimento obbligatorio}
\begin{lstlisting}[style=javacs]
class Prenotazione {
    private final Passeggero passeggero;

    Prenotazione(Passeggero p) {
        if (p == null)
            throw new IllegalArgumentException();
        passeggero = p;
    }
}
class Passeggero {
    private String nome;
}
\end{lstlisting}
\tcblower
\begin{tikzpicture}
  \node[classe1=2.0cm] (r) at (0,0) {Prenotazione};
  \umlclass[2.2cm]{p}{(4.9,-0.12)}{Passeggero}{- nome: String}{}
  \draw[navig] (r.east) -- (r.east -| p.west);
  \estO{(r.east -| p.west)}{1}{passeggero}
\end{tikzpicture}
\regola{Il campo è \code{final}, assegnato nel costruttore e mai \code{null}: molteplicità \code{1}. Non può esistere
una prenotazione senza passeggero, e il costruttore è il punto in cui il codice lo garantisce.}
\end{caso}

\begin{caso}{duecampi}{Due campi dello stesso tipo}
\begin{lstlisting}[style=javacs]
class Volo {
    private final Aeroporto partenza;
    private final Aeroporto arrivo;

    Volo(Aeroporto da, Aeroporto a) {
        partenza = da;
        arrivo = a;
    }
}
class Aeroporto {
    private String codice;
}
\end{lstlisting}
\tcblower
\begin{tikzpicture}
  \node[classe1=1.4cm, minimum height=1.1cm] (v) at (0,0) {Volo};
  \umlclass[2.2cm]{a}{(4.6,0)}{Aeroporto}{- codice: String}{}
  \draw[navig] ($(v.east)+(0,0.3)$) -- node[molt, above, pos=0.85]{1} node[ruolo, above, pos=0.35]{partenza} ($(a.west)+(0,0.3)$);
  \draw[navig] ($(v.east)+(0,-0.3)$) -- node[molt, below, pos=0.85]{1} node[ruolo, below, pos=0.35]{arrivo} ($(a.west)+(0,-0.3)$);
\end{tikzpicture}
\regola{\textbf{Due} campi, \textbf{due} associazioni, ciascuna con il suo ruolo. Non una sola linea con
molteplicità \code{2}: i due aeroporti hanno significati diversi.}
\end{caso}

\begin{caso}{unouno}{Uno a uno, nei due versi}
\begin{lstlisting}[style=javacs]
class Studente {
    private final Badge badge;
    Studente() {
        badge = new Badge(this);   // crea il badge
    }
}
class Badge {
    private final Studente studente;
    Badge(Studente s) {        // riceve lo studente
        studente = s;
    }
}
\end{lstlisting}
\tcblower
\begin{tikzpicture}
  \node[classe1=1.6cm] (s) at (0,0) {Studente};
  \node[classe1=1.6cm] (b) at (4.4,0) {Badge};
  \draw[assocl] (s) -- node[molt, above, pos=0.12]{1} node[ruolo, below, pos=0.16]{studente}
     node[molt, above, pos=0.88]{1} node[ruolo, below, pos=0.86]{badge} (b);
\end{tikzpicture}
\regola{Un campo in \textbf{entrambe} le classi: una sola linea, \textbf{senza frecce}, con un ruolo e una molteplicità
a ogni estremo. \code{1}--\code{1}: un lato crea l'altro e gli passa \code{this}. Se ognuno creasse l'altro
(\code{new Badge()} che fa \code{new Studente()}\ldots) la ricorsione non finirebbe mai (\code{StackOverflowError}).}
\end{caso}

\begin{caso}{unomolti}{Uno a molti, nei due versi}
\begin{lstlisting}[style=javacs]
class Persona {
    private final List<Automobile> automobili =
            new ArrayList<>();
    Persona() {               // almeno un'auto
        automobili.add(new Automobile(this));
    }
    void compraNuova() {
        automobili.add(new Automobile(this));
    }
}
class Automobile {
    private final Persona proprietario;
    Automobile(Persona p) { proprietario = p; }
}
\end{lstlisting}
\tcblower
\begin{tikzpicture}
  \node[classe1=1.6cm] (p) at (0,0) {Persona};
  \node[classe1=1.8cm] (a) at (5.2,0) {Automobile};
  \draw[assocl] (p) -- node[above, font=\scriptsize, yshift=4pt]{Possiede $\blacktriangleright$} (a);
  \estE{(p.east)}{1}{proprietario}
  \estO{(a.west)}{1..*}{automobili}
\end{tikzpicture}
\regola{La \textbf{lista} sta dal lato «uno» (\code{Persona}), il riferimento singolo dal lato «molti»
(\code{Automobile}). \code{1..*}: il costruttore di \code{Persona} crea subito un'automobile. \code{1}: ogni
automobile riceve il proprietario nel costruttore e non lo cambia più.}
\end{caso}

\begin{caso}{moltimolti}{Molti a molti}
\begin{lstlisting}[style=javacs]
class Studente {
    private final List<Esame> esami =
            new ArrayList<>();
    void supera(Esame e) {
        esami.add(e);
        e.registra(this);   // allinea l'altro lato
    }
}
class Esame {
    private final List<Studente> studenti =
            new ArrayList<>();
    void registra(Studente s) { studenti.add(s); }
}
\end{lstlisting}
\tcblower
\begin{tikzpicture}
  \node[classe1=1.6cm] (s) at (0,0) {Studente};
  \node[classe1=1.6cm] (e) at (4.4,0) {Esame};
  \draw[assocl] (s) -- node[molt, above, pos=0.1]{*} node[ruolo, below, pos=0.18]{studenti}
     node[molt, above, pos=0.9]{*} node[ruolo, below, pos=0.84]{esami} (e);
\end{tikzpicture}
\regola{Una collezione in \textbf{ciascuna} classe: \code{*} da entrambe le parti, senza frecce. I metodi devono tenere
allineati i due lati. Se serve un dato sulla coppia (il voto) ci vuole una classe associativa
(\vv{associativacodice}).}
\end{caso}

\begin{caso}{riflessivacodice}{Un campo del tipo della classe stessa}
\begin{lstlisting}[style=javacs]
class Persona {
    private String nome;
    private Persona coniuge;                 // 0..1
    private final List<Persona> figli =
            new ArrayList<>();               // *

    void sposa(Persona p) {
        coniuge = p;
        p.coniuge = this;
    }
    void nasce(Persona f) { figli.add(f); }
}
\end{lstlisting}
\tcblower
\begin{tikzpicture}
  \umlclass[2.6cm]{p}{(0,0)}{Persona}{- nome: String}{}
  % coniuge: nei due versi, a destra
  \draw[assocl] ($(p.east)+(0,0.15)$) -- ++(1.0,0) -- ++(0,1.15) -| ($(p.north)+(0.7,0)$);
  \node[molt, anchor=north west] at ($(p.east)+(0.08,0.08)$) {0..1};
  \node[molt, anchor=south east] at ($(p.north)+(0.64,0.06)$) {0..1};
  \node[ruolo, anchor=south] at ($(p.east)+(-0.1,1.32)$) {coniuge};
  % figli: solo dal genitore, a sinistra
  \draw[navig] ($(p.north)+(-0.7,0)$) -- ++(0,0.95) -| ($(p.west)+(-1.0,0.15)$) -- ($(p.west)+(0,0.15)$);
  \node[molt, anchor=south east] at ($(p.west)+(-0.1,0.22)$) {*};
  \node[ruolo, anchor=north east] at ($(p.west)+(-0.1,0.08)$) {figli};
\end{tikzpicture}
\regola{Un'associazione dalla classe a se stessa (\textbf{riflessiva}). \code{coniuge} va nei due versi
(\code{sposa} aggiorna entrambi): linea senza frecce, \code{0..1} da ogni lato. \code{figli} si naviga solo dal
genitore: freccia, \code{*}.}
\end{caso}

\begin{caso}{rifinterfaccia}{Un campo di tipo interfaccia}
\begin{lstlisting}[style=javacs]
interface Notificatore {
    void invia(String testo);
}
class Email implements Notificatore {
    public void invia(String testo) { /* ... */ }
}
class Ordine {
    private final Notificatore notificatore;
    Ordine(Notificatore n) { notificatore = n; }
    void conferma() {
        notificatore.invia("Ordine confermato");
    }
}
\end{lstlisting}
\tcblower
\begin{tikzpicture}
  \umlclass[2.4cm]{o}{(0,0)}{Ordine}{}{+ conferma()}
  \umlclass[3.0cm]{n}{(5.2,0)}{\stereotipo{interface}{Notificatore}}{}{+ invia(testo: String)}
  \node[classe1=1.5cm] (e) at (5.2,-2.0) {Email};
  \draw[navig] (o.east) -- (n.west);
  \estO{(n.west)}{1}{notificatore}
  \draw[realizza] (e.north) -- (n.south);
\end{tikzpicture}
\regola{La freccia va all'\textbf{interfaccia}, il tipo del campo, non alla classe concreta: \code{Ordine} non sa
nulla di \code{Email}. L'oggetto concreto arriva dal costruttore. È la forma della Dependency Injection e dello
Strategy (\vv{iniezione}).}
\end{caso}

\begin{caso}{rifenum}{Un campo di tipo enumerazione}
\begin{lstlisting}[style=javacs]
enum Stato { BOZZA, PUBBLICATO, ARCHIVIATO }

class Articolo {
    private String titolo;
    private Stato stato = Stato.BOZZA;
    void pubblica() { stato = Stato.PUBBLICATO; }
}
\end{lstlisting}
\tcblower
\begin{tikzpicture}
  \umlclass[3.4cm]{a}{(0,0)}{Articolo}{- titolo: String\\- stato: Stato = BOZZA}{+ pubblica()}
  \node[rectangle split, rectangle split parts=2, draw=q9!75!black, fill=yellow!10, line width=0.7pt,
        font=\scriptsize, text width=2.1cm, align=left, inner sep=2.5pt] (s) at (4.2,0)
        {\centering{\tiny«enumeration»}\\\textbf{Stato}\par\nodepart{two}BOZZA\\PUBBLICATO\\ARCHIVIATO};
  \draw[assocl] (a.east) -- (s.west |- a.east);
\end{tikzpicture}
\regola{Il campo di tipo \code{enum} resta un \textbf{attributo}, con il default (\code{= BOZZA}); l'enumerazione si
disegna a parte con i suoi letterali. La linea fra le due è facoltativa, di solito senza ruoli né molteplicità.}
\end{caso}
